\documentclass[12pt]{article}
\usepackage{url}
\usepackage{layout}

\textwidth=450pt
\textheight=600pt
\oddsidemargin=2pt
\hoffset=0pt
\voffset=0pt

\title{Calculation of phonon spectrum within NEMO-3D code}
\author{Olga L. Lazarenkova}
\date{Last update -- \today}
\begin{document}

\maketitle

\begin{abstract}
This note containes a brief description of  the model, input file  and
referes to some examples of calculation of phonon spectrum.
\end{abstract}

\section{Model}
The current version of the code is dealing with two-parameter Keating
model for calculation of the phonon spectrum. It may be applicable to
covalent semiconductors (C, Si, Ge, etc.)  with 10\% accuracy in the
whole Brillouin zone, except the lack in reproducing flattering of
transversal acoustical dispersion near the Brillouin zone
boundary. The model does not include Coulomb interaction, which leads
to splitting of optical phonon states in the center of the Brillouin
zone in bulk A$_3$B$_5$ semiconductors. Thus phonon dispersion,
computed for A$_3$B$_5$ material systems, may serve as a preliminary
estimation only.

\subsection{Valence-force-field Keating model}
The silmplest way to deal with crystall lattice on atomic level is to
model it by a set of small balls of certain mass, corresponding to
different atoms, connected with springs that represent atomic bonds.

Two-parameter Keating model consider only nearest-neighbour atoms
interactions. The original Keating \cite{Keating_VFFM} formula that
was rewritten later by Martin \cite{Martin_ZnS_VFFM} for the
distorsion energy per unit cell in covalent semiconductors is the
following
\begin{equation}
\label{E_cell}
E_{cell}=\frac{1}{2}\alpha\left({3\over4d^2}\right)\sum_{i=1}^4\left[\Delta\left({\bf r}_i^1\cdot{\bf r}_i^1\right)\right]^2+\frac{1}{2}\sum_{s=1}^2\beta^s\left({3\over4d^2}\right)\sum_{i=1}^4\sum_{j>i}\left[\Delta\left({\bf r}_i^s\cdot{\bf r}_j^s\right)\right]^2.
\end{equation}
Here $\alpha_{ij}$ and $\beta_{ij}$ represent the force constants for bond
length and bond angle distortions in the crystal, respectively;  ${\bf
d}$ and ${\bf r}$ are the equilibrium and real bond length,
respectively; $\Delta\left({\bf r}_i\cdot{\bf r}_j\right)$ is the
change of the dot product of the two bond vectors that start at the
common atom, point along the bond directions, and end at the
\textit{first-neighbor} atoms. 

\subsection{Phonon spectrum}
As a result of thermal fluctuations every $i$-th atom is displaced
from its equilibrium position ${\bf d}_i$ by an amount ${\bf u}({\bf
d}_i)$. The total kinetic energy of the lattice becomes
\begin{equation}
\label{kin}
T=\frac{1}{2}\sum_i\sum_{\xi=x,y,z}M_i\left({\mathrm{d}\: u_\xi({\bf
  d}_i) \over \mathrm{d}\:t}\right)^2.
\end{equation}
The total potential energy of the crystal is assumed to be some
function of the instantaneous positions of all atoms. Within the
\emph{harmonic approximation} the corresponding equation of motion has
the following form
\begin{equation}
\label{eq_motion}
M_i{\mathrm{d}^2\: u_\xi({\bf d}_i) \over \mathrm{d}\:
  t^2}=-\sum_{j}\sum_{\zeta=x,y,z}{\partial^2 E_T \over \partial
  u_\xi({\bf d}_i)\partial u_\zeta({\bf d}_j)} u_\zeta({\bf d}_j).
\end{equation}
Here $E_T$ is the total distortion energy given by summation of
Eq.~(\ref{E_cell}) over all unit cells. Assuming the harmonicity of vibrations, the
solution of the equation of motion (\ref{eq_motion}) may be
represented as
\begin{equation}
\label{U_cluster}
{\bf U}({\bf d}_i,\omega)=\sqrt{M_i}{\bf u}({\bf
d}_i)\exp\left(-i\omega t\right).
\end{equation}
Thus the elements of \emph{dynamical matrix} to be solved for
eigenvalues corresponding to $\omega^2$ have the following form
\begin{equation}
\label{D_cluster}
D_{\xi \zeta}({\bf d}_i,{\bf d}_j)=\left(M_i
  M_j\right)^{-\frac{1}{2}}{\partial^2 E_T \over \partial u_\xi({\bf
  d}_i)\partial u_\zeta({\bf d}_j)}
\end{equation}
The eigenvalues of the dynamical matrix in NEMO-3D are always searched
by PARPACK.


\section{Input file}
In order to be able to compute phonon dispersion you should specify in
{\bf Execution parameters} tree both

-{\it Calculate phonon dispersion}\\
and

-{\it Calculate strain distribution}.\\
The later option is absolutely necessary, since otherwise you are
out of the minimum of the strain energy and the results of phonon
calculations would will be wrong.

Then you should go to {\bf Phonon calculation} subtree of {\bf
Execution parameters} tree and specify

-{\it Phonon model}={\bf VFFM} for {\bf V}alence {\bf F}ield {\bf
 F}orce {\bf M}odel, since it is the only model implemented so far.

-{\it Minimal q$_{ph}$ [nm$^{-1}$]}\\
and

-{\it Maximal q$_{ph}$ [nm$^{-1}$]}\\ 
are three-component vectors of
initial and final phonon wave vector. The phonon spectrum will be
computed in the equidistant points in the direction from minimal to
maximal q$_{ph}$ in the reciprocal space.

-{\it Number of q$_{ph}$ points} is the number of points in which you
 want the phonon spectrum to be computed. For rough estimations of the
 dispersion 10 points should be enough. In order to acheive the better
 smoothness of phonon branches you will have to trade between
 computational time, which scales linearly, and your desire to see nice
 pictures.

-{\it Part of the phonon spectrum} should be set to {\bf Lower} if you
 are interested in {\it acoustic} and {\bf Higher} if you are
 interested in {\it optical} phonons.

-{\it Number of phonon branches} is recommended to be larger then 3
 for better convergency. Also note that the number of phonon branches
 has a physical upper limit, which is three times larger then the
 number of atoms in the solution domain. However, due to the specifics
 of eigensolver, we would not recommend to ask for more then a half of
 the existing eigenstates for tyny systems. You may compute lower and
 upper parts of it in separate runs if you are interested in the whole
 phonon spectrum. For larger system the computational time increases
 dramatically if you are too curious and would like to have too many
 states at once. The recommended upper limit for the number of phonon
 branches is about a dozen.

-{\bf Parpack} sub-subtree would be better left untouched...

-{\bf Output} sub-subtree has only one implemented option so far: {\it
 Phonon dispersion}. If you do not specify it, you only will be able
 to see the corresponding output on the screen after all computations
 completed. If you DO specify it, the file {\tt input\_file\_name.PhonDisp} will
 be rewritten after calculation of energy spectrum completed in every
 requested point of the reciprocal space.

\section{Output file}
The output information about computed phonon dispersion is  saved in 
{\tt input\_file\_name.PhonDisp} file if you have specified it in the
{\bf Output} sub-subtree of {\bf Phonon calculation} subtree of {\bf
Execution parameters} tree.

This is an example of the output file {\tt GeSi\_pa.PhonDisp}  for 8 lowest phonon branches of
a three-dimensional array of pyramidal Ge quantum dots embedded into Si.
{\tiny
\begin{verbatim}
%8 branches of phonon dispersion:

% #   q_x,nm^-1     q_y,nm^-1    q_z,nm^-1  |    q,nm^-1  |   0-br   1-br   2-br   3-br   4-br   5-br   6-br   7-br  

 0  0  0  0  0  0  6.61827e-10  6.73645e-10  0.00168103  0.00173119  0.00173152  0.00177565  0.00178688 
 1  0  0  0.0253684  0.0253684  0.000220557  0.000254217  0.000282209  0.00169869  0.00174823  0.00174912  0.00178307  0.0017895 
 2  0  0  0.0507368  0.0507368  0.000432254  0.000491335  0.00053944  0.00174125  0.00177668  0.00178374  0.00180485  0.00180649 
 3  0  0  0.0761053  0.0761053  0.000627543  0.000700047  0.000756606  0.00177192  0.00178773  0.00182736  0.00182759  0.00184651 
 4  0  0  0.101474  0.101474  0.000801149  0.000876681  0.000932241  0.00178976  0.00179998  0.00184964  0.00188272  0.00191156  
 5  0  0  0.126842  0.126842  0.000950565  0.00102279  0.00107355  0.00181356  0.00181919  0.00187484  0.00193757  0.00198882  
 6  0  0  0.152211  0.152211  0.00107613  0.00114194  0.00118913  0.00184726  0.00184734  0.00190359  0.00199477  0.002059  
 7  0  0  0.177579  0.177579  0.00118052  0.00123821  0.00128562  0.00188414  0.00189174  0.00193679  0.00205174  0.00208816  
 8  0  0  0.202947  0.202947  0.00126752  0.00131607  0.00136814  0.00192735  0.00194541  0.00197588  0.00210548  0.00211059  
 9  0  0  0.228316  0.228316  0.00134084  0.00138027  0.00144115  0.00197271  0.00200557  0.00202243  0.00213049  0.00215376  
 10  0  0  0.253684  0.253684  0.00140389  0.00143526  0.00150824  0.00201758  0.00206897  0.00207569  0.00214967  0.00219585  
 11  0  0  0.279053  0.279053  0.00145979  0.00148471  0.00157163  0.00206223  0.00213145  0.00213257  0.00216921  0.00223249  
 12  0  0  0.304421  0.304421  0.00151121  0.00153118  0.0016323  0.00210751  0.00218063  0.0021925  0.00219272  0.00226477  
 13  0  0  0.329789  0.329789  0.00156006  0.00157627  0.00169053  0.00215365  0.00220529  0.00222256  0.00223252  0.00228604  
 14  0  0  0.355158  0.355158  0.00160744  0.0016207  0.00174629  0.00219862  0.00220773  0.00220952  0.00223209  0.00231815  
 15  0  0  0.380526  0.380526  0.00165352  0.00166438  0.00179924  0.00216534  0.00216878  0.00224668  0.00225257  0.00227263  
 16  0  0  0.405895  0.405895  0.00169748  0.00170626  0.00184854  0.00211923  0.00212568  0.00222733  0.00227091  0.00229348  
 17  0  0  0.431263  0.431263  0.00173684  0.00174375  0.00189214  0.00207715  0.00208555  0.0021863  0.00228861  0.00233422  
 18  0  0  0.456632  0.456632  0.00176627  0.00177161  0.00192503  0.00204574  0.00205564  0.00215483  0.00230115  0.00236492  
 19  0  0  0.482  0.482  0.00177734  0.00178202  0.00193768  0.00203399  0.00204451  0.00214261  0.00230563  0.00237651  
\end{verbatim}
} The first column counts points in the reciprocal space; phonon wave
vector $x,y,z$ components in nm$^{-1}$ are written in the next 3
columns; the 5-th column contains the amplitude of the phonon wave
vector in nm$^{-1}$; and the energy of phonon modes in eV is written in the
rest 8 columns.

The format of the ASCII {\tt input\_file\_name.PhonDisp} file is
designed to allow to plot phonon dispersion curves in matlab. In order
to plot the first two dispersion curves from the output file {\tt
GeSi\_pa.PhonDisp} with blue dots connected with lines you should type
the following commands in matlab:
\begin{verbatim}
>> e=load('GeSi_pa.PhonDisp');
>> figure
>> plot(e(:,5),e(:,6),'b.-');
>> hold on
>> plot(e(:,5),e(:,7),'b.-');
\end{verbatim}

This output file is rewritten after calculation of the energy spectrum is
completed in every requested point of the reciprocal space. Thus if
you see that some rows of the file are filled with zeros, it means
that the calculations have not been completed for this point yet.

\section{Examples}
\subsection{Dispersion of bulk materials}
The bulk material may be built from translation in all 3 directions of
one unit cell. That is you should specify only one shape with
$a_0 \leq L_x=L_y=L_z < 2a_0$, where $a_0$ is unstrained lattice constant of
the material.

{\tt Si.xml} is the example of the input file to compute 12 lower
branches of the phonon dispersion in [001] crystallographic
direction of bulk Si. You may change the {\it Part of the phonon spectrum} to
{\it Higher} and together with the {\it Lower} one it would complete the
whole phonon spectrum of bulk Si.

Note that since the actual unit cell consists of 4 elementary cells,
i.e., it containes 8 atoms instead of 2, the computed spectrum corresponds to the
twice folded real one (see the output file {\tt
  Si.PhonDisp}) and there are 24 eigenstates instead of 6 phonon
branches. You should ``unfold'' it manually.

\subsection{Phonon dispersion in nanostructures}
The calculation of phonon spectrum may be done for all types of
nanostructures generally implemented in NEMO-3D code: quantum wells,
nanowires, quantum dots and their regimented 1D, 2D, and 3D
arrays. Also, you may specify more complicated structures using
combination of different shapes.

{\tt GeSi\_pa.xml} is the example of the input file to compute 12
lower branches of the phonon dispersion in [001] crystallographic
direction of a three-dimensionally regimented array of pyramidal Ge
quantum dots in Si with atomically smooth interfaces. The base of the
dot $L_x=L_y=6.518$nm (i.e.,24 monolayers)and its heigth
$L_z=3.26$nm(i.e.,12 monolayers) are specified in Shape 2. The
periods $D_x=D_y=13.04$nm (i.e.,48 monolayers), and $D_z=6.518$nm
(i.e.,24 monolayers) are specified in Shape 1. The fragment of the
corresponding output file {\tt GeSi\_pa.PhonDisp} was shown in the
previous section.

\section{Conclusion}
The phonon part of the NEMO-3D is still under development... ;)


\begin{thebibliography}{16}

\bibitem{Cardona} 
Peter Y.\ Yu and Manuel Cardona,
\emph{Fundamentals of Semiconductors: Physics and Material Properties} 
3rd, revised and enlarged edition, 
(Springer,N-Y), 2001.

\bibitem{Keating_VFFM}
P.\ N.\ Keating,
Effect of Invariance requirements on the elastic strain energy of
crystals with application to the diamond structure,
\emph{Phys.\ Rev.\ }{\bf 145},2, 637--645 (1966).

\bibitem{Martin_VFFM}
Richard M.\ Martin,
Dielectric screening Model for Lattice Vibrations of Diamond-Structure
Crystals,
\emph{Phys.\ Rev.\ }{\bf 186},3, 871--884 (1969).

\bibitem{Martin_ZnS_VFFM}
Richard M.\ Martin,
Elastic Properties of ZnS Structure Semiconductors,
\emph{Phys.\ Rev.\ B}{\bf 1},10, 4005--4011 (1970).

\bibitem{Pryor}
C.\ Pryor, J.\ Kim, L.\ W.\ Wang, A.\ J.\ Williamson, and A.\ Zunger,
Comparison of two methods for describing the strain profiles in
quantum dots,
\emph{J.\ of Appl.\ Phys.\ }{\bf 83}, 5, 2548--2554 (1998).

\bibitem{NEMO3D}
Gerhard Klimeck, Fabiano Oyafuso, Timothy B.\ Boikin, R.\ Cris Bowen,
and Paul von Allmen,
Development of a Nanoelectronic 3-D (NEMO 3-D) Simulator for
multimillion Atom Simulations and its Application to Alloyed Quantum
Dots,
\emph{CMES}{\bf 3} 5, 601-642 (2002).

\bibitem{Alloys}
An-Ban Chen and Arden Sher,
\emph{Semiconductor alloys: physics and materials engineering}
(Plenum Press, N-Y) 1995.

\bibitem{Harrison_book}
Walter A.\ Harrison,
\emph{Elementary Electronic Structure}
(World Scientific Publishing Co.\ Pte.\ Ltd., London) 1999.

\bibitem{Landolt22}
Landolt-B\"{o}rnstein, Series III, Vol.~22 ({\em Semiconductors\/}), Subvolume
a. \emph{Intrinsic properties of group IV elements, III-V, II-VI, and
  I-VII compounds} (Springer, Berlin, Heidelberg 1987).

\bibitem{Boer}
Karl W.\ B\"{o}er,
\emph{Survey of Semiconductor Physics}, 2nd revised edition (John
Wiley \& Sons, Inc.: New York) 2002.

\bibitem{Maradudin}
A.A~Maradudin,E.W.~Montroll, G.H.~Weiss, and I.P.~Ipatova
\emph{Theory of Lattice Dynamics in the Harmonic Approximation}
(Academic Press, New York and Londod) 1971.

\bibitem{Ewald}
P.P.~Ewald,
\emph{Ann.\ Physik} {\bf 64}, 253 (1921).

\bibitem{Kane}
E.O.\ Kane,
Phonon spectra of diamond and zinc-blende semiconductors,
\emph{Phys.\ Rev.\ B} {\bf 31}, 7865 (1985).

\bibitem{Tersoff}
J.\ Tersoff,
Modeling solid-state chemistry: Interatomic potentials for
multicomponent systems
\emph{Phys.\ Rev.\ B} {\bf 39}, 5566 (1989).

\bibitem{QD_emp_piezo} 
M.A.\ Migliorato, A.G.\ Cullis, M.\ Fearn, and J.H.\ Jefferson,
Atomistic simulation of In$_x$Ga$_{1-x}$As/GaAs
quantum dots with nonuniform composition,
\emph{Physica E} {\bf 13}, 1147 (2002).

\bibitem{Zunger}
H.\ Fu, V.\ Ozoli\c{n}\v{s}, and A.\ Zunger,
Phonons in GaP quantum dots,
\emph{Phys.\ Rev.\ B} {\bf 59}, 2881 (1999).

\end{thebibliography}

\end{document}
